\section{Thảo luận}
Triển vọng của các khung hoạch định dựa trên Quy hoạch lồi hỗn hợp số nguyên (Mixed-Integer Convex Programming - MICP) nằm ở khả năng hội tụ ưu điểm từ hai trường phái chủ đạo đã được phân tích. Cụ thể, phương pháp này hướng tới việc kế thừa tính đầy đủ của các thuật toán dựa trên lấy mẫu (SBP) trong việc khám phá không gian cấu hình phức tạp, đồng thời duy trì năng lực xử lý chính xác các ràng buộc động học và động lực học (kinodynamic constraints) vốn là thế mạnh của các phương pháp tối ưu hóa quỹ đạo (DTO). Lợi ích cốt lõi của việc áp dụng MICP là khả năng đảm bảo tính tối ưu toàn cục (global optimality) trong một khuôn khổ toán học thống nhất, loại bỏ sự tách biệt giữa giai đoạn hoạch định thô và tinh chỉnh quỹ đạo thường thấy trong các chiến lược phân rã truyền thống. Mặc dù sở hữu những ưu điểm lý thuyết vượt trội, sự phổ biến của các kỹ thuật MICP trong thực tế bị hạn chế đáng kể bởi chi phí tính toán và thời gian thực thi. Thực trạng hiện nay cho thấy ngay cả đối với các bài toán có quy mô nhỏ, các thuật toán MICP điển hình có thể yêu cầu thời gian xử lý lên đến vài phút để thiết lập một quỹ đạo duy nhất, gây khó khăn cho các ứng dụng đòi hỏi phản hồi thời gian thực. Gần đây, một số nghiên cứu đã đề xuất các bộ lập kế hoạch không va chạm hoàn toàn dựa trên tối ưu hóa lồi, tuy nhiên phạm vi ứng dụng của chúng vẫn bị giới hạn trong việc lập kế hoạch đường đi thuần túy hình học trong các không gian trạng thái có số chiều thấp. Điều này tạo ra một khoảng trống công nghệ trong việc giải quyết các hệ thống robot có bậc tự do cao với các ràng buộc vi phân phức tạp. Trong bài báo này, chúng tôi tập trung giải quyết một lớp bài toán lập kế hoạch chuyển động đặc thù nhưng đóng vai trò quan trọng đối với các hệ thống có ràng buộc vi phân liên tục. Chúng tôi trình bày một bộ lập kế hoạch mới dựa trên nền tảng MICP nhưng được tối ưu hóa để vượt qua các rào cản về tốc độ thực thi. Điểm khác biệt cốt lõi của nghiên cứu này nằm ở khả năng giải quyết một cách đáng tin cậy các bài toán trong không gian cấu hình số chiều rất cao chỉ trong thời gian tính toán tính bằng giây. Kết quả này đạt được thông qua việc cấu trúc lại bài toán tối ưu hóa thành một chương trình lồi duy nhất, cho phép hệ thống duy trì tính khả thi về mặt động lực học trong khi vẫn đảm bảo hiệu suất vận hành cần thiết cho các môi trường thực tế.
